\documentclass[11pt,a4paper]{article}
\usepackage[margin=25mm]{geometry}
\usepackage{amsmath,amssymb,amsthm,booktabs,hyperref}
\setlength{\emergencystretch}{2em}
\newtheorem{theorem}{Theorem}
\newtheorem{lemma}[theorem]{Lemma}
\newcommand{\Aut}{\operatorname{Aut}}
\newcommand{\gen}[1]{\langle#1\rangle}
\title{Generating pairs from characters and subgroup incidence}
\author{Human Compression project}
\date{}
\begin{document}\maketitle
\begin{abstract}
We give a common exact calculation for generating pairs whose product is
an involution. One subgroup-incidence matrix serves all prescribed element
orders. A complementary stabilizer construction produces orbit representatives
and tests simultaneous inversion. For the O'Nan group we verify two class
obstructions used by Connor and Leemans. The numerical census of generating
pairs is not inferred from the character table alone and remains outside
this component.
\end{abstract}

\section{The objects and the retained character identity}
For a finite group $H$ and integers $p,q>1$, write
\[
 \begin{gathered}
 X_H(p,q)=\{(R,S)\in H^2:|R|=p,\ |S|=q,\ |RS|=2\},\\
 E_H(p,q)=\{(R,S)\in X_H(p,q):\gen{R,S}=H\}.
 \end{gathered}
\]
Put $x_H(p,q)=|X_H(p,q)|$ and $e_H(p,q)=|E_H(p,q)|$.
Pairs are ordered. Restricting a table to $p\leq q$ chooses one order for
the types; it does not further identify a pair with its dual when $p=q$.
Any such additional quotient requires a separate orbit calculation.

We retain the ordinary character theorem of Frobenius, in the form quoted
by Connor--Leemans~\cite{CL}, Theorem~2.1. If $C_i$ are the conjugacy
classes of $H$, with representatives $c_i$, then the number of triples
in $C_i\times C_j\times C_k$ with product one is
\[
 \frac{|C_i||C_j||C_k|}{|H|}
 \sum_{\chi\in\operatorname{Irr}(H)}
 \frac{\chi(c_i)\chi(c_j)\chi(c_k)}{\chi(1)}.
\]
For each character and order define
\[
 v_{\chi,s}=\sum_{|c_i|=s}|C_i|\chi(c_i).
\]
Summing the retained identity over the three order classes gives
\begin{equation}\label{char}
 x_H(p,q)=\frac1{|H|}\sum_\chi
       \frac{v_{\chi,2}v_{\chi,p}v_{\chi,q}}{\chi(1)}.
\end{equation}
Indeed $(R,S)$ corresponds bijectively to $((RS)^{-1},R,S)$ with product
one, and $(RS)^{-1}=RS$ here. All quantities are computed exactly in the
character field; the resulting nonnegative integer is checked, not rounded.
The vectors $v_{\chi,s}$ can be formed once for all desired $(p,q)$.

\section{One incidence calculation for every type}
Choose representatives $H_1,\ldots,H_m$ of every conjugacy class of
subgroups of a fixed finite group $G$, ordered by nondecreasing size.
For each $i,j$ let
\[
 M_{ij}=|\{K\leq H_i:K\text{ is conjugate in }G\text{ to }H_j\}|.
\]
This is a precise weighted subgroup input. A list of subgroup names or
unweighted containment relations alone does not supply $M$.

\begin{theorem}\label{inc}
The matrix $M$ is unit lower triangular. For every pair of orders,
\[
 x_{H_i}(p,q)=\sum_j M_{ij}e_{H_j}(p,q),\qquad
 e_{H_i}(p,q)=x_{H_i}(p,q)-\sum_{j<i}M_{ij}e_{H_j}(p,q).
\]
Consequently the same integral triangular calculation recovers the whole
vector of generating-pair counts from the corresponding all-pair vectors.
\end{theorem}
\begin{proof}
Every pair in $X_{H_i}(p,q)$ generates a unique subgroup $K\leq H_i$.
Its contribution belongs to exactly one of the sets $E_K(p,q)$. Conjugation
preserves both orders and the product relation, and gives a bijection
$E_K(p,q)\simeq E_{H_j}(p,q)$ whenever $K$ is conjugate to $H_j$.
Partitioning by the actual generated subgroup proves the first formula.
If $K\leq H_i$ and $|K|\geq|H_i|$, then equality of orders forces
$K=H_i$. There is precisely one such subgroup in its own conjugacy class
and none in a different equal-order class. Thus $M$ is unit lower triangular,
and moving the preceding terms gives the second formula. The same proof
applies in every coordinate of a vector of types; $M$ is independent of
those types.
\end{proof}

\begin{lemma}\label{free}
The action of $\Aut(H)$ on $E_H(p,q)$ is free. Its number of orbits is
$e_H(p,q)/|\Aut(H)|$.
\end{lemma}
\begin{proof}
An automorphism fixing both generators fixes every word in them and hence
every element of $H$. Every stabilizer is therefore trivial, so every orbit
has the indicated group order. This freeness need not hold on $X_H(p,q)$;
dividing the all-pair count by $|\Aut(H)|$ is not the required orbit count.
\end{proof}

Theorem~\ref{inc} is Hall's subgroup inversion arranged as one common
calculation. It is not a new inversion theorem. Formula~\eqref{char} and
Lemma~\ref{free} likewise retain the methods attributed in~\cite{CL}.
The source already reuses subgroup data in its algorithm; the vector
organization supplies an explicit shared interface, not a claim that its
authors recomputed every subgroup for each arithmetic operation.

\section{Constructing the representatives and testing inversion}
Let $A\leq\Aut(G)$ be a specified automorphism group, and let $J$ be the
set of nonidentity involutions in $G$. The map
\[
 (R,S)\longmapsto(S,RS),\qquad (s,t)\longmapsto(ts^{-1},s)
\]
is an $A$-equivariant bijection between product-involution pairs and
$G\times J$. Moreover $\gen{R,S}=\gen{s,t}$.

\begin{theorem}\label{orbits}
Choose one representative $s$ from every $A$-orbit of elements of order
$q$. For each such $s$, choose representatives for the point stabilizer
$A_s$ on
\[
 J_{s,p}=\{t\in J:|ts^{-1}|=p\}.
\]
Retain exactly those with $\gen{s,t}=G$. The resulting pairs
$(ts^{-1},s)$ represent each $A$-orbit in $E_G(p,q)$ exactly once.
\end{theorem}
\begin{proof}
Every pair has its second coordinate moved to a chosen $s$ by an element
of $A$. Its new involution coordinate belongs to $J_{s,p}$. This set is
$A_s$-invariant, because its definition uses only $s$, multiplication and
element order. A further element of $A_s$ moves that coordinate to a chosen
representative. Generation is also invariant, proving coverage after the
final test. If two retained pairs are equivalent, their second coordinates
first lie in the same chosen orbit and are therefore equal. The transporting
automorphism then lies in $A_s$, so their chosen involution coordinates
are equal. This proves uniqueness.
\end{proof}

This is the same stabilizer argument used in the acquired L071 enumeration
proof, applied to a different equivariant set. It does not require IP of a
tuple or that $s$ be an involution. Using only $\gen s$-orbits on $J$, as
in the source's constructive route, still covers the candidates but requires
its subsequent full-isomorphism removal. The full stabilizer avoids that
particular redundancy. Neither construction may use an unproved expected
count as its stopping certificate.

For a generating pair with $p,q>2$, at most one $\theta\in\Aut(G)$
can simultaneously send $(R,S)$ to $(R^{-1},S^{-1})$. Indeed the ratio of
two such automorphisms fixes the generating pair. If it exists then
$\theta^2=1$ and $\theta\ne1$. The simultaneous inversion criterion
for regular maps is retained from~\cite{CL}, Section~5: no such
automorphism gives the chiral case, inner inversion gives the non-orientable
case for the centreless simple group, and outer inversion gives the
extension case. Testing a pair's two elements separately for invertibility
is only a necessary condition, since the same automorphism must do both.

\section{The two O'Nan obstructions}
Let $G$ be the O'Nan simple group, and acquire its named ordinary character
table and that of $\Aut(G)=G:2$ from the Atlas data in CTblLib.
We also acquire the stored fusion from $G$-classes to $G:2$-classes.
These finite-group identifications and tables are retained inputs, not
identified solely from the displayed orders.

\begin{lemma}\label{inverse}
Let $G\trianglelefteq A$ have index two. Let $i$ be a conjugacy class of
$G$, write $\iota(i)$ for its inverse class and $f(i)$ for its class in
$A$. An element of class $i$ has an inner inverter precisely when
$\iota(i)=i$. It has an inverter in $A\setminus G$ precisely when either
\[
 \iota(i)\ne i,\quad f(\iota(i))=f(i),
 \quad\text{or}\quad
 \iota(i)=i,\quad |f^{-1}(f(i))|=1.
\]
\end{lemma}
\begin{proof}
The first assertion is the definition of conjugacy to the inverse in $G$.
Every $A$-class contained in $G$ is the union of either one or two
$G$-classes, since $G$ has index two. If the inverse class is different
but fuses with $i$, an inverter exists in $A$ and none exists in $G$,
so it is outside $G$. If the inverse class is $i$ and the fusion fibre
has two classes, every outer conjugation exchanges them and cannot invert
the element. If the fibre has one class, an outer conjugate can be brought
back to the element by inner conjugation, giving an outer centralizer
element. Composing it with an inner inverter gives an outer inverter.
These cases exhaust the possibilities.
\end{proof}

The exact table calculation gives the following data; the last columns
count classes whose elements admit the specified inverter.
\begin{center}\begin{tabular}{rrrr}\toprule
order&$G$-classes&inner inverter&outer inverter\\\midrule
8&2&2&0\\16&4&4&0\\20&2&0&2\\31&2&0&2\\\bottomrule
\end{tabular}\end{center}
Hence no generating pair with an element of order 20 or 31 admits inner
inversion; no pair with an element of order 8 or 16 admits outer inversion.
This recovers exactly Lemmas 5.2 and 5.3 of~\cite{CL}, using the retained
simultaneous-inversion interface. It does not prove any positive type or
the remaining isolated exclusions in Theorem 1.1.

\section{Verification and the remaining input}
The exact programs compare character counts with all literal pairs in
$S_3,A_4,S_4,A_5$. A separate control forms the complete conjugate-subgroup
incidence matrices for $S_3,D_8,A_4,S_4,A_5$, performs the vector inversion,
and compares every resulting count against direct subgroup generation.
These finite controls test the implementation; the uniform results are the
preceding proofs, not an extrapolation from the examples.

For O'Nan, the program computes all character counts and the complete
inverse/fusion data. They are only all-pair counts. The full weighted
subgroup incidence and subgroup character inputs required for
Theorem~\ref{inc}, or a completed representative enumeration under
Theorem~\ref{orbits}, have not been supplied here. In particular this
component does not certify the source's 796772 generating-pair orbits,
its 136 numerical table entries or its full chirality spectrum. No
whole-original admission, new conjecture, Lean statement or compression
measurement is assigned to it.
\begin{thebibliography}{9}
\bibitem{CL} T. Connor and D. Leemans, Algorithmic enumeration of regular
maps, \emph{Ars Math. Contemp.}10 (2016),211--222.
\bibitem{Hall} P. Hall, The Eulerian functions of a group,
\emph{Quart. J. Math. Oxford}7 (1936),134--151.
\bibitem{Frobenius} G. Frobenius, Ueber Gruppencharaktere,
\emph{Sitzungsber. Preuss. Akad. Wiss. Berlin} (1896),985--1021.
\bibitem{Atlas} J.H. Conway, R.T. Curtis, S.P. Norton, R.A. Parker and
R.A. Wilson, \emph{Atlas of Finite Groups}, Oxford University Press,1985.
\end{thebibliography}
\end{document}
